%%%%%%%%%%%%%%%%%%%%%%%%%%%%%%%%%%%%%%%%%%%%%%%%%%%%%%%%%%%%%%%%%%%
\section{Workload 2: Whole-Vocabulary Readout}
\label{sec:LargeK}
%%%%%%%%%%%%%%%%%%%%%%%%%%%%%%%%%%%%%%%%%%%%%%%%%%%%%%%%%%%%%%%%%%%
% (1) what the method is and why its K is not a choice; (2) our in-house
% reproduction at 8B, as a watchlist; (3) the cost of full readout;
% (4) what would make it tractable, and which of those we have shown.

This section evaluates the mappings on the second workload:
whole-vocabulary readout, where the library is set by the model's
vocabulary rather than by anyone's choices -- the
$K\approx1.3\times10^{5}$ rung of Fig.~\ref{fig:ladder}. We describe
the published method and why its $K$ is structural, report an
in-house reproduction at watchlist scale with its measured fidelity
at \cam precision, and price what full readout would cost a
\cam-resident design against streaming the same bank from HBM.
\inote[aiviolet]{Aug 31: section intro added per your ``frame this
section like V'' note}

\subsection{A workload whose $K$ nobody selected}
\label{subsec:Jspace}
The grids in Sec.~\ref{sec:Results} are representative of the
$K\approx10^{2}$--$10^{3}$ regime, where the library grows with the
number of behaviors somebody thinks to watch. A recently published
readout method sets $K$ a different way. It reports that a model holds words ``in
mind'' without emitting them, and that for every token in the vocabulary
there is an internal pattern meaning that word is currently active,
readable whether or not the word is ever
said~\cite{anthropic2026workspace}. In the spider example of
Sec.~\ref{sec:Introduction}, the word appears in neither prompt nor
answer, and replacing the internal \emph{spider} pattern with the one
for \emph{ant} changes the answer from eight to six. The same work
reports that the readout catches misbehavior beyond the staged safety
test: ``manipulation'' was active while a model produced falsified
data during a pre-release audit.

For our purposes the relevant property is structural rather than
behavioral. Reading this space means comparing the internal state against
one stored pattern \emph{per vocabulary token} -- roughly 128K signatures
for a Llama-family tokenizer -- at several layers, every token. Nobody
chose that library size; it falls out of the vocabulary. And it is
plausibly a floor, since the published method tracks single words, and
phrases would only enlarge the library.

\subsection{An in-house reproduction, at watchlist scale}
\label{subsec:Watchlist}
We reproduced the readout in-house on Llama-3.1-8B: a layer-by-layer
ranked list of what the model is currently representing. About 13 of 32
layers carry a usable readout at this scale, which is the source of the
``13--25 usable layers'' factor used in Sec.~\ref{sec:Scaling}. Each
signature is precomputed offline, so a watched word costs one comparison
per token, and we stored 47 deception-adjacent signatures per monitored
layer as \cam rows.

In a staged scenario in which the model is told it is about to be
replaced and is given a personnel file it was not meant to use, ten watch
words were tracked by exact full-vocabulary rank across the 2{,}048
prompt tokens the model reads before generating anything. Five of them --
``scandal,'' ``evidence,'' ``leaked,'' ``fake'' and ``shutdown'' --
reached rank~0 at interpretable moments in the reading, ``threat''
reached rank~1 and ``blackmail'' rank~2, and none of the ten appears in
the prompt text.\inote[aiviolet]{corrected Aug 21: the v01 draft said
three words at rank 0, which is how the number appears in every
downstream deck. The 07-16-26 meeting notes are the source of record and
say five, plus threat at 1 and blackmail at 2} The reply the model then produced was cooperative and
unremarkable, i.e., not something a text classifier would score as
harmful. Ground truth for the reproduction is the model's own
full-precision readout, so no labeled deception data is required
anywhere.

These are preliminary numbers: one model, one staged scenario, ten watch
words in the demonstration and 47 stored signatures.\footnote{Several of
the fidelity and scaling values in this subsection were read from result
charts rather than from result files; verification against the team's
files is queued, and the scaling sweep re-runs with three known fixes
before any of it is reused externally.}


\subsection{Fidelity at \cam precision}
\label{subsec:JspaceFidelity}
The deception grid of Sec.~\ref{sec:Results} needs labeled activations.
Readout does not: the full-precision, full-dimension dot product against
the stored patterns \emph{is} the ground truth, so a compressed
implementation can be scored directly against the exact readout it is
trying to reproduce. Compression can therefore be measured here without
labels, and the deception numbers of Sec.~\ref{sec:Results} are not
assumed to carry over to it.

The setup takes J-space vectors at a fixed layer, gathers residual-stream
activations from Wikitext, records which stored vectors fire above a
global threshold under the exact computation, and asks how well each
\cam-compatible encoding reproduces that list. Two readouts are scored:
\emph{threshold}, where every row that fires is reported and the metric
is per-row firing F1, and \emph{top-1}, where the single closest row is
reported and the metric is agreement with the exact readout's choice.
The bank is 47 deception-adjacent vectors per layer, as in
Sec.~\ref{subsec:Watchlist}.

Fig.~\ref{fig:jspace}(a,b) plots both against storage bits per \cam
row, which is the axis that connects fidelity to
Sec.~\ref{sec:Hardware}: bits per row set the bank size, and the bank
size sets what has to be held or moved. An uncompressed fp16 row at
$d=8{,}192$ is 131{,}072 bits and reproduces the exact readout by
definition. Three observations.

\textbf{Two orders of magnitude of compression cost less than it might:}
A 3-bit Euclidean code at $\sim$$10^{3}$ bits per row -- a $128\times$
reduction -- reaches \hbox{0.96} top-1 agreement and 0.81 per-row firing
F1. Keeping the projection in full precision and quantizing nothing but
the width (the exact-$L_2$ series) reaches 1.00 top-1 at 2{,}048 bits and
0.84 firing F1 at 4{,}096, so most of the remaining gap at $10^{3}$ bits
is the quantizer rather than the projection.

\textbf{The top-1 readout is markedly cheaper than the threshold
readout:} At every code width the top-1 curve sits above the firing-F1
curve, by 0.15--0.2 in the range that matters. Reporting \emph{which row
is closest} survives compression better than reporting \emph{which rows
cross a threshold}, because the first is a comparison among rows and the
second is a comparison against an absolute value that quantization
shifts. For a \cam this is a matchline-mode result: best-match mode
tolerates a cheaper cell than threshold mode does, at equal fidelity.

\textbf{Hamming and thermometer codes do not reach the same operating
point here:} Both rise with code width to roughly 0.74 top-1 and 0.49
firing F1 at $10^{3}$ bits, and the Hamming series then falls sharply at
4{,}096 bits, where the encoding switches to a direct sign hash without a
learned projection. That is the opposite ordering from
Sec.~\ref{sec:Results}, where at a matched 384 bits the sign code was
ahead. The two workloads are asking different questions -- separating two
labeled classes there, reproducing a specific ranked list here -- and the
metric a readout needs appears to be the one a multi-bit cell provides.
Sec.~\ref{subsec:Devices} takes up what this would mean for the array
if it survives the checks queued in
Sec.~\ref{subsec:NextMeasurements}.
\inote[aiviolet]{v03: per the v02 note, the design-space consequence of
this disagreement now has its own paragraph in Sec.~VII-D, and the
projection-artifact check is item (3) of the measurement plan in
Sec.~VII-E}

\begin{figure*}[t]
\centering
\includegraphics[width=7.0in]{fig_jspace_v1.pdf}
\caption{Whole-vocabulary readout reproduced through a \cam.
(a) Per-row firing F1 and (b) top-1 agreement with the exact J-lens,
against storage bits per \cam row, for a bank of 47 vectors per layer;
the star marks an uncompressed fp16 row (131{,}072 bits), which
reproduces the exact readout by definition. (c) Per-row macro F1 as the
library grows from $10^{2}$ to $10^{4}$ tracked vectors at 124$+$4
cells per row. Values digitized from the source deck; see the footnote
in Sec.~\ref{subsec:JspaceFidelity} for the limitations that apply to
panel~(c).}
\label{fig:jspace}
\end{figure*}

\subsection{Fidelity as the library grows}
\label{subsec:JspaceScaling}
Compression fidelity at a 47-row bank does not establish anything about a
$10^{5}$-row one, so the same pipeline was swept over library size.
Fig.~\ref{fig:jspace}(c) reports per-row macro F1 from $K=10^{2}$ to
$K=10^{4}$.\footnote{Panels (a) and (b) are digitized from the charts in
the group's 07-30-26 result deck; the underlying result files were not
available. Panel (c) is from the same deck and carries three limitations
its author states: an epsilon-inside-square-root bug drove the $K=32$ and
$K=320$ points to zero and they are dropped here; mean-token pooling was
used by mistake, so the residual-stream vectors differ slightly from the
per-token vectors a deployment would see; and \cam top-$k$ clipping is
assumed away, i.e., a row that fires is taken to be selected. A re-run
with all three fixed is queued.}

Three things follow. \textbf{The library stays alive:} the fraction of
rows that remain distinguishable is essentially 1.0 out to $10^{4}$
tracked vectors for a trained projection, with a fixed SVD projection
beginning to shed rows at about $10^{4}$ (to 0.88). Nothing collapses.
\textbf{Per-row fidelity nonetheless degrades:} macro F1 falls from 0.755
at $K=10^{2}$ to 0.43 at $K=10^{4}$ for the 3-bit trained projection.
Against an always-fire baseline that drops from 0.10 to 0.005 over the
same range, 0.43 is still nearly two orders of magnitude better than
trivial, but it is not a fidelity anybody would deploy against.
\textbf{The degradation is not a quantization artifact:} the
full-precision projection degrades faster, from 0.69 to 0.30, and is
overtaken by the 3-bit trained projection above $K\approx10^{3}$. What
gets harder as the library grows is the discrimination, not the
arithmetic, which is a result in favor of the compressed
implementation rather than against it.

In engineering terms, holding $10^{4}$ rows is not the problem and
reading them faithfully is. That places the open work on the encoding
and the readout mode rather than on the array, and it is the reason
Sec.~\ref{sec:Hardware} prices capacity and traffic separately.

\subsection{What full readout would cost}
\label{subsec:ReadoutCost}
Priced at 70B geometry, an uncompressed bank of 128K signatures at
$d=8{,}192$, fp16, is \qty{2.1}{\giga\byte} per layer. Streaming that on
every token at 1{,}000 tokens/s over 25 monitored layers is roughly
\qty{53}{\tera\byte\per\second}, which is not a workload anyone runs. The
useful part of the answer, however, stays small: only on the order of 25
patterns are meaningfully active at any one time, so the operation is
detection rather than readout -- scan everything, return the few that
fired -- which is the question a threshold-mode \cam answers natively.

Three things would make it tractable, in increasing order of what they
require.

\textbf{A watchlist:} Store the few hundred words already known to be
worth watching, which drops $K$ by three orders of magnitude. Our 8B
demonstration in Sec.~\ref{subsec:Watchlist} is this case. A watchlist
catches only what somebody thought to list, and the words that matter
most may be the ones nobody anticipated, so it is a first rung toward
full readout rather than a substitute for it.

\textbf{Threshold matching instead of ranking:} A \cam can report which
rows are within $\tau$ without scoring the rest, and with roughly 25
active patterns out of 128K nearly every query returns nothing.

\textbf{Compressing the stored side:} At the code widths
Sec.~\ref{sec:Results} measured, 384--2{,}048 bits per row against
131{,}072 bits for an uncompressed fp16 row, the whole 128K bank is
6--\qty{33}{\mega\byte} per layer (Fig.~\ref{fig:capacity}), which fits
on a die. The requirement is that it not move. Only the first of the three has been
demonstrated here; the other two are design arguments.

\begin{figure}[t]
\centering
\includegraphics[width=3.4in]{fig_capacity_v1.pdf}
\caption{Bank capacity per monitored layer against library size, for
uncompressed fp16 rows and for the two code widths measured in
Sec.~\ref{sec:Results}, with two on-chip capacity references. A
whole-vocabulary bank at the measured code widths is
6--\qty{33}{\mega\byte} per layer; the same bank uncompressed is
\qty{2.1}{\giga\byte}.}
\label{fig:capacity}
\end{figure}
